\documentclass[11pt,a4paper]{article}
\usepackage[margin=2.5cm]{geometry}
\usepackage{amsmath,amssymb,amsthm}
\usepackage{tikz}
\usepackage{booktabs}
\usepackage{array}
\usepackage{hyperref}

\usetikzlibrary{positioning, calc}

\theoremstyle{plain}
\newtheorem{theorem}{Theorem}
\newtheorem{proposition}[theorem]{Proposition}
\newtheorem{corollary}[theorem]{Corollary}
\newtheorem{lemma}[theorem]{Lemma}

\theoremstyle{definition}
\newtheorem{definition}[theorem]{Definition}
\newtheorem{example}[theorem]{Example}

\theoremstyle{remark}
\newtheorem{remark}[theorem]{Remark}

\newcommand{\Z}{\mathbb{Z}}
\newcommand{\R}{\mathbb{R}}
\newcommand{\F}{\mathbb{F}}
\newcommand{\Aut}{\operatorname{Aut}}
\newcommand{\Spec}{\operatorname{Spec}}
\newcommand{\tr}{\operatorname{tr}}

\title{The Spectrum of the Hypercube:\\
  Cartesian Products, Boolean Fourier Analysis, and Expansion}
\author{Robin Langer}
\date{}

\begin{document}
\maketitle

\begin{abstract}
We give a self-contained account of the adjacency and Laplacian spectra of the
$n$-dimensional hypercube graph~$Q_n$, derived from the Cartesian product
decomposition $Q_n = K_2^{\square n}$. The eigenvalues $n - 2k$ with
multiplicities~$\binom{n}{k}$ arise naturally from Boolean Fourier analysis on
$\Z_2^n$, connecting graph spectrum to harmonic analysis on the Boolean cube.
We determine the exact threshold ($n = 7$) at which hypercubes cease to be
Ramanujan, characterise their expansion properties, and relate the spectral
structure to coding theory via the Hamming scheme.
\end{abstract}


% ==========================================================================
\section{Introduction}
% ==========================================================================

The $n$-dimensional hypercube $Q_n$ is among the most studied objects in
combinatorics. Its vertex set is $\Z_2^n = \{0,1\}^n$, with two vertices
adjacent when they differ in exactly one coordinate. As a Cayley graph of the
elementary abelian 2-group $\Z_2^n$ with generators $\{e_1, \ldots, e_n\}$
(the standard basis vectors), it inherits a rich algebraic structure that
completely determines its spectrum.

The spectrum of $Q_n$ is well known, but its derivation from three distinct
viewpoints --- Cartesian products, character theory, and Boolean Fourier
analysis --- reveals connections that are individually elementary yet
collectively illuminating. This note develops all three approaches, then
applies them to the Ramanujan property and expansion.


% ==========================================================================
\section{The hypercube as Cartesian product}
% ==========================================================================

\begin{definition}
The \emph{Cartesian product} $G \mathbin{\square} H$ of graphs $G$ and $H$
has vertex set $V(G) \times V(H)$, with $(g_1, h_1) \sim (g_2, h_2)$ iff
either $g_1 = g_2$ and $h_1 \sim h_2$ in $H$, or $h_1 = h_2$ and
$g_1 \sim g_2$ in $G$.
\end{definition}

The complete graph $K_2$ has two vertices and one edge, with adjacency
matrix $A_{K_2} = \bigl(\begin{smallmatrix} 0 & 1 \\ 1 & 0 \end{smallmatrix}\bigr)$
and spectrum $\{+1, -1\}$. The hypercube is the iterated Cartesian product:
\[
  Q_n \;=\; \underbrace{K_2 \mathbin{\square} K_2 \mathbin{\square}
    \cdots \mathbin{\square} K_2}_{n \text{ copies}}.
\]

\begin{theorem}[Cartesian product spectral theorem]
\label{thm:cart-spec}
If $G$ has eigenvalues $\lambda_1, \ldots, \lambda_m$ and $H$ has
eigenvalues $\mu_1, \ldots, \mu_k$, then $G \mathbin{\square} H$ has
eigenvalues $\{\lambda_i + \mu_j : 1 \le i \le m,\; 1 \le j \le k\}$.
\end{theorem}

\begin{proof}
The adjacency matrix of $G \mathbin{\square} H$ decomposes as
\[
  A_{G \square H} = A_G \otimes I_k + I_m \otimes A_H.
\]
If $A_G v = \lambda v$ and $A_H w = \mu w$, then
$A_{G \square H}(v \otimes w) = (\lambda + \mu)(v \otimes w)$.
Since the tensor products of eigenbases form a basis for $\R^{mk}$,
this accounts for all eigenvalues.
\end{proof}

\begin{corollary}[Spectrum of $Q_n$]
\label{cor:hypercube-spec}
The adjacency eigenvalues of $Q_n$ are
\[
  \lambda_k = n - 2k, \qquad
  \text{with multiplicity } \binom{n}{k}, \qquad
  k = 0, 1, \ldots, n.
\]
\end{corollary}

\begin{proof}
By induction on Theorem~\ref{thm:cart-spec}. Each eigenvalue of $Q_n$ is a
sum $\varepsilon_1 + \cdots + \varepsilon_n$ where each
$\varepsilon_i \in \{+1, -1\}$. If exactly $k$ of the $\varepsilon_i$ equal
$-1$, the sum is $(n-k) - k = n - 2k$, and there are $\binom{n}{k}$ such
choices. As a consistency check: $\sum_{k=0}^n \binom{n}{k} = 2^n = |V(Q_n)|$.
\end{proof}


% ==========================================================================
\section{Character-theoretic derivation}
% ==========================================================================

Since $Q_n$ is a Cayley graph of the abelian group $\Z_2^n$, its eigenvalues
are given by the characters of $\Z_2^n$ evaluated on the connection set.

The irreducible characters of $\Z_2^n$ are indexed by subsets
$S \subseteq [n] = \{1, \ldots, n\}$ (equivalently, by elements of the
Pontryagin dual $\widehat{\Z_2^n} \cong \Z_2^n$):
\[
  \chi_S(x) = (-1)^{\sum_{i \in S} x_i}, \qquad x \in \Z_2^n.
\]
The connection set of $Q_n$ is $C = \{e_1, \ldots, e_n\}$, so the eigenvalue
corresponding to $\chi_S$ is:
\[
  \lambda_S = \sum_{j=1}^n \chi_S(e_j)
  = \sum_{j=1}^n (-1)^{[j \in S]}
  = (n - |S|) - |S|
  = n - 2|S|.
\]
The eigenvalue $n - 2k$ has multiplicity equal to the number of subsets of
size $k$, namely $\binom{n}{k}$, recovering Corollary~\ref{cor:hypercube-spec}.


% ==========================================================================
\section{Boolean Fourier analysis}
% ==========================================================================

The character-theoretic approach is a special case of Fourier analysis on
$\Z_2^n$. Any function $f : \Z_2^n \to \R$ can be expanded as
\[
  f(x) = \sum_{S \subseteq [n]} \hat{f}(S)\, \chi_S(x),
\]
where $\hat{f}(S) = \mathbb{E}_{x}[f(x)\,\chi_S(x)]$ and expectation is
over the uniform distribution. The characters $\{\chi_S\}$ form an
orthonormal basis under the inner product
$\langle f, g \rangle = \mathbb{E}_x[f(x)g(x)]$.

The adjacency operator $A$ acts on functions by
$(Af)(x) = \sum_{y \sim x} f(y)$. Since $\chi_S$ is an eigenfunction of $A$
with eigenvalue $\lambda_S = n - 2|S|$, the Fourier decomposition
simultaneously diagonalises $A$. The \emph{level sets}
\[
  W_k = \operatorname{span}\{\chi_S : |S| = k\}
\]
are the eigenspaces of $Q_n$, and $\dim W_k = \binom{n}{k}$.

\begin{remark}
The level sets $W_k$ are precisely the \emph{Reed--Muller codes} of order
$k$ modulo order $k{-}1$. The connection between hypercube spectrum and coding
theory runs deeper: the eigenvalues of $Q_n$ are the Krawtchouk polynomial
values $K_k(j; n) = \sum_{s=0}^k (-1)^s \binom{j}{s}\binom{n-j}{k-s}$
evaluated at $j = 1$, and these polynomials define the Hamming association
scheme.
\end{remark}


% ==========================================================================
\section{Laplacian spectrum}
% ==========================================================================

For an $n$-regular graph, the Laplacian is $L = nI - A$, so the Laplacian
eigenvalues of $Q_n$ are $\mu_k = n - \lambda_k = 2k$ with multiplicity
$\binom{n}{k}$.

\begin{proposition}
The algebraic connectivity (Fiedler value) of $Q_n$ is $\mu_1 = 2$, with
multiplicity $n$, for all $n \ge 1$.
\end{proposition}

\begin{proof}
The smallest nonzero Laplacian eigenvalue is $\mu_1 = 2\cdot 1 = 2$,
attained by the $n$ characters $\chi_{\{i\}}$ for $i = 1, \ldots, n$.
\end{proof}

\begin{remark}
The Fiedler eigenvectors $\chi_{\{i\}}(x) = (-1)^{x_i}$ are precisely the
coordinate functions. Each one is constant on the two halves
$\{x : x_i = 0\}$ and $\{x : x_i = 1\}$, reflecting the fact that a single
coordinate bisection (removing $2^{n-1}$ edges) is the minimum bisection of
$Q_n$. The constant Fiedler value $\mu_1 = 2$ means the ``bottleneck'' of a
hypercube is always a single-coordinate cut, regardless of dimension.
\end{remark}


% ==========================================================================
\section{The Ramanujan threshold}
% ==========================================================================

\begin{definition}
A connected $d$-regular graph $G$ is \emph{Ramanujan} if every eigenvalue
$\lambda$ of $A$ with $|\lambda| < d$ satisfies $|\lambda| \le 2\sqrt{d-1}$.
\end{definition}

For $Q_n$, the second-largest eigenvalue is $\lambda_2 = n - 2$, so $Q_n$ is
Ramanujan iff
\begin{equation}\label{eq:ram}
  n - 2 \;\le\; 2\sqrt{n-1}.
\end{equation}

\begin{proposition}
The hypercube $Q_n$ is Ramanujan if and only if $n \le 6$.
\end{proposition}

\begin{proof}
Squaring both sides of \eqref{eq:ram} (valid since both sides are
nonneg\-ative for $n \ge 2$):
\[
  (n-2)^2 \le 4(n-1)
  \;\Longleftrightarrow\;
  n^2 - 4n + 4 \le 4n - 4
  \;\Longleftrightarrow\;
  n^2 - 8n + 8 \le 0.
\]
The roots are $n = 4 \pm 2\sqrt{2} \approx 1.17$ and $6.83$, so the
inequality holds for $n \le 6$ (checking $n = 1$ separately: $\lambda_2 = -1
\le 0$). At $n = 7$: $\lambda_2 = 5 > 2\sqrt{6} \approx 4.899$.
\end{proof}

\begin{remark}
The Ramanujan property is an asymptotic spectral measure of
expansion quality. That $Q_n$ fails at $n = 7$ does not mean $Q_7$ is a poor
expander --- only that its second eigenvalue grows linearly ($\lambda_2 =
n - 2$) while the Ramanujan bound grows as $O(\sqrt{n})$. As a family,
$\{Q_n\}$ is not a family of Ramanujan graphs, but it is an
\emph{expander family}: the normalised spectral gap
$\gamma_n = (\lambda_1 - \lambda_2)/\lambda_1 = 2/n$ is bounded below
by $2/n > 0$, and the Cheeger constant $h(Q_n) = 1$ is optimal
among vertex-transitive graphs.
\end{remark}


% ==========================================================================
\section{Expansion and isoperimetry}
% ==========================================================================

The constant spectral gap of $Q_n$ has direct consequences for expansion.

\begin{theorem}[Cheeger inequality for $Q_n$]
The edge-isoperimetric number (Cheeger constant) of $Q_n$ satisfies:
\[
  1 = \frac{\mu_1}{2} \;\le\; h(Q_n) \;\le\; \sqrt{2n \cdot \mu_1}
  = 2\sqrt{n}.
\]
In fact, $h(Q_n) = 1$ exactly: for any subset $S \subseteq V(Q_n)$ with
$|S| \le 2^{n-1}$, the number of edges between $S$ and its complement
satisfies $|\partial S| \ge |S|$.
\end{theorem}

This is a consequence of Harper's vertex-isoperimetric theorem, which
states that among all sets of size $m$ in $Q_n$, the initial segment in
binary order (a Hamming ball) minimises the boundary. The edge version
follows from the Lindsey--Loomis--Whitney inequality.

\begin{proposition}[Mixing time]
The lazy random walk on $Q_n$ (stay with probability $1/2$, move to a
random neighbour with probability $1/2$) has mixing time
$\Theta(n \log n)$.
\end{proposition}

\begin{proof}[Proof sketch]
Each coordinate is an independent lazy random walk on $K_2$, which mixes in
$\Theta(\log n)$ steps (coupon collector for all $n$ coordinates to have
been touched). The spectral gap of the lazy walk is $\mu_1 / (2n) = 1/n$,
giving $t_{\mathrm{mix}} = \Theta(n \log n)$ by the standard gap bound.
\end{proof}


% ==========================================================================
\section{The Hamming scheme}
% ==========================================================================

The hypercube sits inside the \emph{Hamming scheme} $H(n, 2)$, the
association scheme on $\Z_2^n$ where the $i$-th relation $R_i$ connects
pairs at Hamming distance $i$. The adjacency matrix of $Q_n$ is the first
relation matrix $A_1$.

The eigenmatrix of $H(n, 2)$ is the matrix of Krawtchouk polynomials:
\[
  P_{kj} = K_k(j; n) = \sum_{s=0}^{k} (-1)^s \binom{j}{s}
  \binom{n-j}{k-s}.
\]
The adjacency eigenvalue of $Q_n$ corresponding to the $k$-th eigenspace is
$P_{k,1} = K_k(1; n) = n - 2k$, which recovers our formula once more.

\begin{remark}
Distance-regularity of $Q_n$ means it generates the full Bose--Mesner
algebra of $H(n, 2)$ through its spectral idempotents.
The spectral decomposition
\[
  A_1 = \sum_{k=0}^n (n-2k)\, E_k, \qquad
  E_k = \frac{1}{2^n} \sum_{j=0}^n K_k(j; n)\, A_j
\]
makes $Q_n$ a microcosm of the entire scheme: knowing the $Q_n$ spectrum is
equivalent to knowing all distance relations.
\end{remark}


% ==========================================================================
\section{Automorphisms and representation theory}
% ==========================================================================

The automorphism group of $Q_n$ is the hyperoctahedral group
$\Aut(Q_n) \cong \Z_2^n \rtimes S_n$ of order $2^n \cdot n!$, acting by
coordinate permutations and reflections (bit flips). This is the same as the
signed permutation group, or equivalently the wreath product $\Z_2 \wr S_n$.

The eigenspaces $W_k = \operatorname{span}\{\chi_S : |S| = k\}$ are
$\Aut(Q_n)$-invariant, since $S_n$ permutes subsets of fixed size and
$\Z_2^n$ acts by $\pm 1$ on each character. As a representation of $S_n$
(the non-translation part of the automorphism group), $W_k$ carries the
$k$-th exterior power of the standard representation:
\[
  W_k \cong \bigwedge^k \R^n \quad \text{as } S_n\text{-modules}.
\]

\begin{proposition}
Each eigenspace $W_k$ is irreducible as an $S_n$-module (corresponding to
the partition $(n-k, 1^k)$, i.e.\ a hook), and is characterised by the Young
diagram with one row of length $n-k$ and a column of height $k$.
\end{proposition}

\begin{remark}
This is the graph-theoretic analogue of Hodge theory: the Laplacian
eigenspaces of $Q_n$ decompose $L^2(\Z_2^n)$ into ``$k$-forms,'' with $W_k$
playing the role of harmonic $k$-forms on the Boolean cube. The formal
similarity is not accidental --- $Q_n$ is the 1-skeleton of the $n$-cube
polytope, and $W_k$ is the image of the $k$-th simplicial cohomology under
the combinatorial Hodge correspondence.
\end{remark}


% ==========================================================================
\section{Spectral data}
% ==========================================================================

\begin{table}[h]
\centering
\caption{Adjacency spectrum and key invariants of $Q_n$ for $1 \le n \le 10$.}
\label{tab:spectra}
\smallskip
\renewcommand{\arraystretch}{1.2}
\small
\begin{tabular}{@{}c r r p{5.6cm} c c c@{}}
\toprule
$Q_n$ & $|V|$ & $\deg$ & Adjacency spectrum ($\lambda^{\text{mult}}$)
  & $\lambda_2$ & Gap & Ram.\ \\
\midrule
$Q_1$ &     2 &  1 & $1^1,\; {-}1^1$
  & $-1$ & 2 & \checkmark \\
$Q_2$ &     4 &  2 & $2^1,\; 0^2,\; {-}2^1$
  & $0$ & 2 & \checkmark \\
$Q_3$ &     8 &  3 & $3^1,\; 1^3,\; {-}1^3,\; {-}3^1$
  & $1$ & 2 & \checkmark \\
$Q_4$ &    16 &  4 & $4^1,\; 2^4,\; 0^6,\; {-}2^4,\; {-}4^1$
  & $2$ & 2 & \checkmark \\
$Q_5$ &    32 &  5 & $5^1,\; 3^5,\; 1^{10},\; {-}1^{10},\; {-}3^5,\; {-}5^1$
  & $3$ & 2 & \checkmark \\
$Q_6$ &    64 &  6 & $6^1,\; 4^6,\; 2^{15},\; 0^{20},\; {-}2^{15},\; {-}4^6,\; {-}6^1$
  & $4$ & 2 & \checkmark \\
$Q_7$ &   128 &  7 & $7^1,\; 5^7,\; 3^{21},\; 1^{35},\; \ldots$ (symmetric)
  & $5$ & 2 & $\times$ \\
$Q_8$ &   256 &  8 & $8^1,\; 6^8,\; 4^{28},\; 2^{56},\; 0^{70},\; \ldots$
  & $6$ & 2 & $\times$ \\
$Q_9$ &   512 &  9 & $9^1,\; 7^9,\; 5^{36},\; 3^{84},\; 1^{126},\; \ldots$
  & $7$ & 2 & $\times$ \\
$Q_{10}$ & 1024 & 10 & $10^1,\; 8^{10},\; 6^{45},\; 4^{120},\; \ldots$
  & $8$ & 2 & $\times$ \\
\bottomrule
\end{tabular}
\end{table}

Several patterns are visible in Table~\ref{tab:spectra}:
\begin{itemize}
\item The spectrum is always symmetric about zero (since $Q_n$ is bipartite).
\item The multiplicities form row $n$ of Pascal's triangle.
\item The spectral gap $\lambda_1 - \lambda_2 = 2$ is constant.
\item The Ramanujan property fails precisely at $n = 7$.
\end{itemize}


% ==========================================================================
\section{Concluding remarks}
% ==========================================================================

The hypercube is that rare object whose spectrum can be derived in at least
four independent ways --- Cartesian products, character theory, Fourier
analysis, and association schemes --- each of which opens a different window
onto its structure. The multiplicities $\binom{n}{k}$ are simultaneously
dimensions of exterior powers, sizes of Hamming spheres, and weights of
Boolean functions of degree~$k$. The constant spectral gap of~2 is
simultaneously a fact about tensor products, about the coupon collector
problem, and about coordinate bisections.

That the Ramanujan threshold falls at $n = 7$ is a pleasant arithmetic
coincidence: $4 + 2\sqrt{2} \approx 6.83$. One might ask whether the
non-Ramanujan hypercubes $Q_n$ for $n \ge 7$ can be ``repaired'' by passing
to a quotient. The halved cube $\frac{1}{2}Q_n$ (identifying antipodal
vertices) is $\binom{n}{2}$-regular with spectrum governed by even-weight
Krawtchouk polynomials, and its Ramanujan status is a subtler question left
for future investigation.

\bigskip

{\raggedright\noindent\textbf{Computational verification.}
All spectral data was verified numerically against direct eigenvalue
computation of the $2^n \times 2^n$ adjacency matrix for $n \le 12$;
see the accompanying \texttt{hypercube\_spectrum.py}.\par}

\end{document}
